\documentclass[10pt,a4paper]{amsart}
\usepackage{amsmath,amssymb,amsthm,graphicx}
% For commutative diagrams we use Xy-pic
\usepackage[all]{xy}

% We use 2cell for 2-commutative diagrams.
\xyoption{2cell}
% Theorem environments.
%
\theoremstyle{plain}
\newtheorem{theorem}[subsection]{Theorem}
\newtheorem{proposition}[subsection]{Proposition}
\newtheorem{lemma}[subsection]{Lemma}

\theoremstyle{definition}
\newtheorem{definition}[subsection]{Definition}
\newtheorem{example}[subsection]{Example}
\newtheorem{exercise}[subsection]{Exercise}
\newtheorem{situation}[subsection]{Situation}

\theoremstyle{remark}
\newtheorem{remark}[subsection]{Remark}
\newtheorem{remarks}[subsection]{Remarks}

\numberwithin{equation}{subsection}

% Macros
%
\def\lim{\mathop{\mathrm{lim}}\nolimits}
\def\colim{\mathop{\mathrm{colim}}\nolimits}
\def\Spec{\mathop{\mathrm{Spec}}}
\def\Hom{\mathop{\mathrm{Hom}}\nolimits}
\def\Ext{\mathop{\mathrm{Ext}}\nolimits}
\def\SheafHom{\mathop{\mathcal{H}\!\mathit{om}}\nolimits}
\def\SheafExt{\mathop{\mathcal{E}\!\mathit{xt}}\nolimits}
\def\Sch{\mathit{Sch}}
\def\Mor{\mathop{\mathrm{Mor}}\nolimits}
\def\Ob{\mathop{\mathrm{Ob}}\nolimits}
\def\Sh{\mathop{\mathit{Sh}}\nolimits}
\def\NL{\mathop{N\!L}\nolimits}
\def\CH{\mathop{\mathrm{CH}}\nolimits}
\def\proetale{{pro\text{-}\acute{e}tale}}
\def\etale{{\acute{e}tale}}
\def\QCoh{\mathit{QCoh}}
\def\Ker{\mathop{\mathrm{Ker}}}
\def\Im{\mathop{\mathrm{Im}}}
\def\Coker{\mathop{\mathrm{Coker}}}
\def\Coim{\mathop{\mathrm{Coim}}}

% Boxtimes
%
\DeclareMathSymbol{\boxtimes}{\mathbin}{AMSa}{"02}

%
% Macros for moduli stacks/spaces
%
\def\QCohstack{\mathcal{QC}\!\mathit{oh}}
\def\Cohstack{\mathcal{C}\!\mathit{oh}}
\def\Spacesstack{\mathcal{S}\!\mathit{paces}}
\def\Quotfunctor{\mathrm{Quot}}
\def\Hilbfunctor{\mathrm{Hilb}}
\def\Curvesstack{\mathcal{C}\!\mathit{urves}}
\def\Polarizedstack{\mathcal{P}\!\mathit{olarized}}
\def\Complexesstack{\mathcal{C}\!\mathit{omplexes}}
% \Pic is the operator that assigns to X its picard group, usage \Pic(X)
% \Picardstack_{X/B} denotes the Picard stack of X over B
% \Picardfunctor_{X/B} denotes the Picard functor of X over B
\def\Pic{\mathop{\mathrm{Pic}}\nolimits}
\def\Picardstack{\mathcal{P}\!\mathit{ic}}
\def\Picardfunctor{\mathrm{Pic}}
\def\Deformationcategory{\mathcal{D}\!\mathit{ef}}
\newenvironment{reference}{}{}
\newenvironment{slogan}{}{}
\newenvironment{history}{}{}
\usepackage{cleveref}
\crefname{originallemma}{original lemma}{original lemmas}
\Crefname{originallemma}{Original lemma}{Original lemmas}
\makeatletter
\newlabel{lemma-non-irreducible-minimal-type-genus-at-least-one}{{0C7B}{source}{}{sourcepointer-1}{}}
\newlabel{lemma-non-irreducible-minimal-type-genus-at-least-one@cref}{{[originallemma][0][]0C7B}{[1][0][]source}{}{}{}}
\makeatother
\crefname{originallemma}{original lemma}{original lemmas}
\Crefname{originallemma}{Original lemma}{Original lemmas}
\makeatletter
\newlabel{lemma-contract-picard-group}{{0C7J}{source}{}{sourcepointer-2}{}}
\newlabel{lemma-contract-picard-group@cref}{{[originallemma][0][]0C7J}{[1][0][]source}{}{}{}}
\makeatother
\crefname{originallemma}{original lemma}{original lemmas}
\Crefname{originallemma}{Original lemma}{Original lemmas}
\makeatletter
\newlabel{lemma-contract}{{0C77}{source}{}{sourcepointer-3}{}}
\newlabel{lemma-contract@cref}{{[originallemma][0][]0C77}{[1][0][]source}{}{}{}}
\makeatother
\title{A numerical type of genus at most zero has Picard group the integers}
\author{The Stacks Project Authors}
\date{stacks-project:models.tex:local-snapshot-20261012; GNU FDL 1.2 or later}
\begin{document}
\maketitle
\noindent\textit{Editorial note.} The counted unit is the Stacks Project Lemma lemma-picard-group-genus-nonpositive at source lines 1221-1225 together with its proof at source lines 1227-1238, a straight adjacent pairing in the chapter file. The chapter is the upstream file models.tex, which is the file name recorded in the source record; the pool directory carries the longer name stacks-models and is not used as the chapter name anywhere in this card.\par\medskip
\noindent\textit{Editorial note.} This chapter's candidate index records the statement body and the proof body without their closing markers, so the retained spans used here are the full environments, exactly as in the earlier cards of this chapter.\par\medskip
\noindent\textit{Editorial note.} The retained proof names three results of the same chapter: the classification excluding minimal types of nonpositive genus, the comparison of the Picard groups under contraction, and the invariance of the genus under contraction. All three are declared as non-load-bearing navigation pointers so that the references resolve; the counted argument is the induction itself and the structural reduction, not the cited classification. Their numbers are the stable Stacks tags taken from the shared tag index, not guessed values.\par\medskip
\noindent\textit{Editorial note.} Neither the retained statement nor the retained proof contains a citation, so no bibliography entry is supplied and the bibliography backend is none.\par\medskip
\noindent\textit{Editorial note.} The retained preamble is the chapter preamble's Xy-pic setup at lines 14-18, the theorem declarations and mathematical macros at lines 154-217, and the three no-op wrapper environments.\par\medskip
\noindent\textit{Allowed standard prerequisites (editorial).}\begin{itemize}
\item Numerical types, their genus and their Picard group. A numerical type is a collection of integers n, m of subscript i, a of subscript i j, w of subscript i and g of subscript i satisfying the symmetry and recurrence conditions of the chapter; the genus of the type is the alternating sum assembled from those integers and its Picard group is the cokernel of the matrix of divided intersection numbers, as in the neighbouring lemmas of the same section. Equivalence of numerical types and the operation of contracting a minus-one index are described in the same section; the author's own notation is retained inside the counted span.. The counted lemma computes the Picard group for all types of genus at most zero, and it is the base of the induction that the following lemmas of the section use.
\end{itemize}\medskip
\ifdefined\hypertarget\hypertarget{sourcepointer-1}{}\fi
\noindent\textit{Source reference.} Original lemma 0C7B: focused-seeds/AGP/sources/stacks-snapshot-20261011/tags.txt, source lines 15326--15326. Navigation pointer invoked inside the retained proof when the author excludes the minimal case; the counted argument is the induction and the contraction step, and the cited classification of minimal types is an allowed input that is not reproduced inside the span. The number is the stable Stacks tag of the target, looked up in the shared tag index.\par\medskip
\ifdefined\hypertarget\hypertarget{sourcepointer-2}{}\fi
\noindent\textit{Source reference.} Original lemma 0C7J: focused-seeds/AGP/sources/stacks-snapshot-20261011/tags.txt, source lines 15334--15334. Navigation pointer invoked inside the retained proof for the comparison of the two Picard groups; the comparison itself is not reproduced inside the span and is taken as an allowed input of the section. The number is the stable Stacks tag of the target, looked up in the shared tag index.\par\medskip
\ifdefined\hypertarget\hypertarget{sourcepointer-3}{}\fi
\noindent\textit{Source reference.} Original lemma 0C77: focused-seeds/AGP/sources/stacks-snapshot-20261011/tags.txt, source lines 15322--15322. Navigation pointer invoked inside the retained proof for the invariance of the genus under the contraction; that invariance is an allowed input and is not reproduced inside the span. The number is the stable Stacks tag of the target, looked up in the shared tag index.\par\medskip
\begin{lemma}
\label{lemma-picard-group-genus-nonpositive}
Let $n, m_i, a_{ij}, w_i, g_i$ be a numerical type $T$.
If the genus $g$ of $T$ is $\leq 0$, then $\Pic(T) = \mathbf{Z}$.
\end{lemma}

\begin{proof}
By induction on $n$. If $n = 1$, then the assertion is clear.
If $n > 1$, then $T$ is not minimal by
Lemma \ref{lemma-non-irreducible-minimal-type-genus-at-least-one}.
After replacing $T$ by an equivalent type
we may assume $n$ is a $(-1)$-index.
By Lemma \ref{lemma-contract-picard-group}
we find $\Pic(T) \subset \Pic(T')$.
By Lemma \ref{lemma-contract} we see that the genus
of $T'$ is equal to the genus of $T$ and we conclude by
induction.
\end{proof}

\end{document}
